% !TeX root = ../main.tex
% Lezione 4.pdf, pp. 6--14. Le prove dei criteri operativi e della
% proprietà duale dell'infimo rendono espliciti i passaggi della fonte.
\section{Intervalli e retta reale estesa}
% Fonte: Lezione 4, pp.~6--8

\begin{definizione}[Intervallo]
Un sottoinsieme non vuoto $I\subseteq\R$ è un \emph{intervallo} se,
ogni volta che $x,y\in I$ e $x<z<y$, anche $z$ appartiene a $I$.
\end{definizione}

In un intervallo si trovano quindi tutti i punti compresi fra due
suoi elementi. Per estremi finiti $a<b$ le principali possibilità sono
\begin{align*}
  [a,b]&=\{x\in\R:a\le x\le b\},&
  (a,b)&=\{x\in\R:a<x<b\},\\
  [a,b)&=\{x\in\R:a\le x<b\},&
  (a,b]&=\{x\in\R:a<x\le b\}.
\end{align*}
I primi due sono detti rispettivamente intervallo \emph{chiuso} e
\emph{aperto}; negli altri due è incluso un solo estremo. Anche
$[a,a]=\{a\}$ è un intervallo.\footnote{La definizione richiede
che un intervallo sia non vuoto. Altre convenzioni ammettono anche
l'insieme vuoto come intervallo.}

Per $I=[a,b]$, i numeri $a,b$ sono gli \emph{estremi}, $b-a$ è
l'\emph{ampiezza}, $(b-a)/2$ il \emph{raggio} e $(a+b)/2$ il
\emph{centro}. Le stesse espressioni individuano ampiezza e centro
anche se uno o entrambi gli estremi non sono inclusi.

Per descrivere intervalli illimitati introduciamo i simboli
$-\infty$ e $+\infty$, ordinati in modo che
\[
  -\infty<x<+\infty\qquad\forall x\in\R.
\]
La \emph{retta reale estesa} è
$\overline\R=\R\cup\{-\infty,+\infty\}$. Le semirette della retta
reale si scrivono, per esempio,
\begin{align*}
  (-\infty,c]&=\{x\in\R:x\le c\},&
  (-\infty,c)&=\{x\in\R:x<c\},\\
  [a,+\infty)&=\{x\in\R:x\ge a\},&
  (a,+\infty)&=\{x\in\R:x>a\}.
\end{align*}
In queste scritture $\pm\infty$ non appartengono agli intervalli,
che sono sottoinsiemi di $\R$; anche tutta $\R=(-\infty,+\infty)$
è un intervallo.

\section{Limitatezza, maggioranti e minoranti}
% Fonte: Lezione 4, pp.~8--10

\begin{definizione}
Un insieme non vuoto $A\subseteq\R$ è \emph{superiormente limitato}
se esiste $M\in\R$ tale che $a\le M$ per ogni $a\in A$.
Un tale numero è un \emph{maggiorante} di $A$.
È \emph{inferiormente limitato} se esiste $m\in\R$ tale che
$m\le a$ per ogni $a\in A$; un tale numero è un \emph{minorante}.
\end{definizione}

L'insieme è \emph{limitato} se lo è sia superiormente sia inferiormente:
\[
  \exists m,M\in\R:\quad m\le a\le M\quad\forall a\in A.
\]
Equivalentemente, esiste $N\ge0$ per cui $|a|\le N$ per ogni
$a\in A$. Infatti, da due estremi $m,M$ basta scegliere
$N=\max\{|m|,|M|\}$; viceversa, $|a|\le N$ equivale a
$-N\le a\le N$. Un insieme che non è limitato si dice
\emph{illimitato}; si può anche specificare la direzione in cui
manca un limite, parlando di insieme superiormente o inferiormente
illimitato.

Maggioranti e minoranti non devono appartenere all'insieme e non
sono determinati in modo univoco. Per esempio, sono possibili scelte
anche molto lontane dagli estremi:
\begin{itemize}
  \item per $A=[0,1)$ si possono prendere $m=-38$ e $M=251$;
  \item per $B=\{1/n:n\in\N^*\}$ si possono prendere
    $m=-1719{,}3$ e $M=16$, perché $0<1/n\le1$;
  \item per $C=\N$, $0$ è un minorante ma non esistono maggioranti,
    come seguirà dal principio di Archimede
    (teorema~\ref{teo:archimede-reali}).
\end{itemize}

Per $A=[0,3)$, in particolare, sono maggioranti $3$, $5$, $11$,
$197$ e, più in generale, tutti e soli i reali $M\ge3$.
Sono minoranti $0$, $-13$, $-\pi$, $-171$ e, in generale, tutti e soli
i reali $m\le0$. La condizione vale per ogni elemento di $A$,
anche se $3\notin A$.

\section{Massimo e minimo}
% Fonte: Lezione 4, p.~11

\begin{definizione}
Un \emph{massimo} di un insieme non vuoto $A\subseteq\R$ è un
maggiorante che appartiene ad $A$; un \emph{minimo} è un minorante
che appartiene ad $A$. Si scrive
\begin{align*}
  M=\max A&\quad\Longleftrightarrow\quad
    M\in A\text{ e }a\le M\ \forall a\in A,\\
  m=\min A&\quad\Longleftrightarrow\quad
    m\in A\text{ e }m\le a\ \forall a\in A.
\end{align*}
\end{definizione}

Se esistono, massimo e minimo sono unici. Per esempio, se $M_1$ e
$M_2$ fossero entrambi massimi, l'appartenenza di ciascuno ad $A$
imporrebbe $M_1\le M_2$ e $M_2\le M_1$, quindi $M_1=M_2$.
Per il minimo si ragiona allo stesso modo.

L'intervallo $[0,3)$ ha minimo $0$ e non ha massimo: per ogni
$a\in[0,3)$ il punto $(a+3)/2$ è ancora nell'intervallo ed è
strettamente maggiore di $a$. Quindi l'essere limitato non assicura
che un insieme possieda un massimo o un minimo.

\section{Estremo superiore ed estremo inferiore}
% Fonte: Lezione 4, p.~12

\begin{definizione}[Supremo e infimo]
Il numero $S\in\R$ è l'\emph{estremo superiore}, o \emph{supremo},
di $A$ se è il più piccolo dei suoi maggioranti. Il numero $s\in\R$
è l'\emph{estremo inferiore}, o \emph{infimo}, se è il più grande
dei suoi minoranti. In formule,
\begin{align*}
  S=\sup A&=\min\{M\in\R:a\le M\ \forall a\in A\},\\
  s=\inf A&=\max\{m\in\R:m\le a\ \forall a\in A\},
\end{align*}
quando tali estremi esistono.
\end{definizione}

Anche questi valori, quando esistono, sono unici. A differenza del
massimo e del minimo, possono non appartenere ad $A$. Se un massimo
esiste, è anche il più piccolo maggiorante: essendo un elemento di
$A$, è minore o uguale a qualunque altro maggiorante. Pertanto
\[
  \max A=\sup A\quad\text{se il massimo esiste},\qquad
  \min A=\inf A\quad\text{se il minimo esiste}.
\]
Viceversa, se $\sup A\in A$ esso è il massimo, e se $\inf A\in A$
esso è il minimo.

\subsection{Il teorema di esistenza degli estremi}
% Fonte: Lezione 4, pp.~13--14

\begin{teorema}[Esistenza del supremo e dell'infimo]
\label{teo:esistenza-estremi}
Ogni sottoinsieme non vuoto di $\R$ superiormente limitato ha
un estremo superiore reale. Ogni sottoinsieme non vuoto di $\R$
inferiormente limitato ha un estremo inferiore reale.
\end{teorema}
\begin{proof}
Sia $A\subseteq\R$ non vuoto e superiormente limitato. L'insieme
dei suoi maggioranti
\[
  \mathcal M_A=\{M\in\R:a\le M\ \forall a\in A\}
\]
è non vuoto. Per definizione, ogni elemento di $A$ è minore o
uguale a ogni elemento di $\mathcal M_A$. Applichiamo l'assioma
di Dedekind ai due insiemi $A$ e $\mathcal M_A$: esiste $S\in\R$
tale che
\[
  a\le S\le M\qquad\forall a\in A,\ \forall M\in\mathcal M_A.
\]
La prima disuguaglianza dice che $S$ è un maggiorante di $A$;
la seconda che è minore o uguale a ogni maggiorante. Dunque
$S=\min\mathcal M_A=\sup A$.

Per l'infimo si può applicare lo stesso ragionamento all'insieme
$-A=\{-a:a\in A\}$. Se $A$ è inferiormente limitato, $-A$ è
superiormente limitato e ha supremo $T$. Cambiando segno alle
disuguaglianze si ottiene che $-T$ è il più grande minorante di
$A$, ossia $\inf A=-\sup(-A)$.
\end{proof}

Il teorema è una conseguenza della completezza. Sostituendo $\R$
con $\Q$ l'enunciato diventa falso: l'insieme
\[
  A=\{x\in\Q:x\ge0,\ x^2\le2\}
\]
è non vuoto e superiormente limitato, ma non ha supremo in $\Q$.
Un eventuale minimo maggiorante razionale sarebbe infatti un
separatore razionale fra gli insiemi $A$ e $B$ già considerati
nell'esempio sull'assioma di Dedekind. Considerando lo stesso
insieme come sottoinsieme di $\R$, il suo supremo è $\sqrt2$;
il suo infimo è $0$, che è anche minimo.

\subsection{Estremi infiniti}

Per gli insiemi non vuoti si estende la notazione ponendo
$\sup A=+\infty$ quando $A$ è superiormente illimitato e
$\inf A=-\infty$ quando è inferiormente illimitato.
Si tratta di valori nella retta estesa, non di massimi o minimi
reali. Per esempio, $\sup\N=+\infty$ e $\inf\N=\min\N=0$.

\section{La caratterizzazione operativa di supremo e infimo}
% Fonte: Lezione 4, p.~14; lezione 5, p.~2

\begin{proposizione}
\label{prop:criterio-estremi}
Sia $A\subseteq\R$ non vuoto. Un numero $S\in\R$ è il suo supremo
se e solo se valgono entrambe le condizioni
\[
  \begin{cases}
    a\le S&\forall a\in A,\\
    \forall\varepsilon>0\ \exists a\in A:\ a>S-\varepsilon.
  \end{cases}
\]
Analogamente, $s\in\R$ è l'infimo se e solo se
\[
  \begin{cases}
    s\le a&\forall a\in A,\\
    \forall\varepsilon>0\ \exists a\in A:\ a<s+\varepsilon.
  \end{cases}
\]
\end{proposizione}
\begin{proof}
Se $S=\sup A$, esso è un maggiorante. Se per qualche
$\varepsilon>0$ non esistesse $a\in A$ con $a>S-\varepsilon$,
tutti gli elementi sarebbero minori o uguali a $S-\varepsilon$:
quest'ultimo sarebbe un maggiorante più piccolo di $S$, assurdo.

Viceversa, supponiamo che le due condizioni valgano. $S$ è un
maggiorante e ogni $M<S$ non può esserlo: scegliendo
$\varepsilon=S-M>0$ si trova $a\in A$ con $a>M$. Quindi $S$ è
il più piccolo dei maggioranti.

Per l'infimo, se $s=\inf A$ e mancassero elementi sotto
$s+\varepsilon$, quest'ultimo sarebbe un minorante più grande.
Viceversa, se le due condizioni per $s$ valgono, ogni $m>s$ viene
escluso dai minoranti scegliendo $\varepsilon=m-s$ e trovando
$a\in A$ con $a<m$. Dunque $s$ è il più grande minorante.
\end{proof}

Il criterio traduce il ruolo di estremo in due verifiche separate:
essere un limite per tutti gli elementi e poter essere avvicinato
arbitrariamente dall'insieme. L'elemento $a$ può dipendere da
$\varepsilon$. Se l'estremo appartiene ad $A$, lo si può scegliere
per ogni $\varepsilon$; il criterio è particolarmente utile quando
l'estremo non è raggiunto.

\begin{schema}
\begin{tikzpicture}[x=1cm,y=1cm]
  \fill[coloreschema!12] (4.3,.52) rectangle (6.8,.78);
  \fill[coloreschema!12] (1.2,-.98) rectangle (3.7,-.72);
  \foreach \quota in {.65,-.85} {
    \draw[asse] (0,\quota) -- (8,\quota);
  }
  \node[anchor=east] at (-.2,.65) {$S=\sup A$};
  \node[anchor=east] at (-.2,-.85) {$s=\inf A$};
  \foreach \pos/\testo in {4.3/{S-\varepsilon},6.8/S} {
    \draw (\pos,.55) -- (\pos,.75);
    \node[below=6pt] at (\pos,.65) {$\testo$};
  }
  \foreach \pos/\testo in {1.2/s,3.7/{s+\varepsilon}} {
    \draw (\pos,-.95) -- (\pos,-.75);
    \node[below=6pt] at (\pos,-.85) {$\testo$};
  }
  \fill[coloreschema] (6.1,.65) circle[radius=2pt];
  \node[above=6pt,coloreschema] at (6.1,.65) {$a\in A$};
  \fill[coloreschema] (2.1,-.85) circle[radius=2pt];
  \node[above=6pt,coloreschema] at (2.1,-.85) {$a\in A$};
\end{tikzpicture}
\par\smallskip
Per ogni $\varepsilon>0$, la fascia $(S-\varepsilon,S]$ deve
contenere almeno un punto di $A$; analogamente per $[s,s+\varepsilon)$.
Il colore evidenzia la fascia di ricerca, senza affermare che tutti
i suoi punti appartengano ad $A$.
\end{schema}
